%%======================================================================
%%  NT-1: Literature Extraction & Convention Fixing
%%        for SCT Nonlocal Form Factor Computation
%%
%%  This document collects formulas from the literature needed to
%%  explicitly compute F_1(\Box/\Lambda^2) and F_2(\Box/\Lambda^2)
%%  from the SCT spectral action Tr(f(D^2/\Lambda^2)).
%%
%%  RESEARCH DOCUMENT: formulas extracted from cited papers.
%%  No original derivations --- gaps noted explicitly.
%%======================================================================
\documentclass[12pt,a4paper]{article}

%% ---- packages -------------------------------------------------------
\usepackage[utf8]{inputenc}
\usepackage[T1]{fontenc}
\usepackage{lmodern}
\usepackage{amsmath,amssymb,amsthm,mathtools}
\usepackage[margin=2.5cm]{geometry}
\usepackage{booktabs}
\usepackage{enumitem}
\usepackage{hyperref}
\usepackage{xcolor}
\usepackage{fancyhdr}
\usepackage{titlesec}
\usepackage{longtable}

%% ---- theorem-like environments --------------------------------------
\newtheorem{proposition}{Proposition}[section]
\newtheorem{lemma}[proposition]{Lemma}
\newtheorem{corollary}[proposition]{Corollary}
\theoremstyle{definition}
\newtheorem{definition}[proposition]{Definition}
\newtheorem{convention}[proposition]{Convention}
\theoremstyle{remark}
\newtheorem{remark}[proposition]{Remark}

%% ---- custom commands ------------------------------------------------
\newcommand{\Tr}{\operatorname{Tr}}
\newcommand{\tr}{\operatorname{tr}}
\newcommand{\Rsq}{R_{\mu\nu\rho\sigma}R^{\mu\nu\rho\sigma}}
\newcommand{\Ricsq}{R_{\mu\nu}R^{\mu\nu}}
\newcommand{\Weylsq}{C_{\mu\nu\rho\sigma}C^{\mu\nu\rho\sigma}}
\newcommand{\dd}{\mathrm{d}}
\newcommand{\ii}{\mathrm{i}}
\newcommand{\Id}{\mathrm{Id}}
\DeclareMathOperator{\sgn}{sgn}

%% ---- annotation commands --------------------------------------------
\newcommand{\fromlit}[1]{\textcolor{blue!70!black}{\textsf{[#1]}}}
\newcommand{\oursign}{\textcolor{teal!80!black}{\textsf{[our convention]}}}
\newcommand{\gap}{\textcolor{red!80!black}{\textsf{[GAP]}}}
\newcommand{\needscheck}{\textcolor{orange!80!black}{\textsf{[needs sign check]}}}
\newcommand{\exactlit}{\textcolor{blue!60!black}{\textsf{[exact from lit.]}}}
\newcommand{\approxlit}{\textcolor{orange!70!black}{\textsf{[approximate]}}}

%% ---- sign correction annotation -------------------------------------
\newcommand{\signcorrection}[1]{%
  \par\noindent\colorbox{red!10}{%
    \parbox{\dimexpr\linewidth-2\fboxsep}{%
      \textcolor{red!50!black}{\textbf{[Sign correction]}}
      \textcolor{red!50!black}{#1}%
    }%
  }\par%
}

%% ---- page style -----------------------------------------------------
\pagestyle{fancy}
\fancyhf{}
\fancyhead[L]{\small\textsc{NT-1: Literature \& Conventions for Form Factors}}
\fancyhead[R]{\small\thepage}
\renewcommand{\headrulewidth}{0.4pt}

%% ---- title ----------------------------------------------------------
\title{\textbf{NT-1: Literature Extraction \& Convention Fixing\\[4pt]
       for SCT Nonlocal Form Factor Computation}}
\author{David Alfyorov}
\date{March 2026 --- Working Document}

%%======================================================================
\begin{document}
\maketitle

\begin{center}
\small\textit{Credits: Aliaksandr Samatyia contributed research-assistance and workflow support.}
\end{center}

\thispagestyle{empty}

\begin{abstract}
This document collects all formulas, conventions, and computational
recipes needed to explicitly derive the nonlocal form factors
$F_1(\Box/\Lambda^2)$ and $F_2(\Box/\Lambda^2)$ from the SCT spectral
action $\Tr(f(D^2/\Lambda^2))$.  All formulas are extracted from the
literature; no original derivations are performed.  Each formula is
annotated with its source, equation number, and whether it matches our
sign conventions.  Gaps are noted explicitly.

\medskip\noindent
\textbf{Key references:}
\begin{itemize}[nosep]
\item Barvinsky--Vilkovisky (BV): Nucl.\ Phys.\ B \textbf{282} (1987) 163;
  B \textbf{333} (1990) 471, 512; arXiv:0911.1168.
\item Avramidi: Phys.\ Lett.\ B \textbf{236} (1990) 443;
  B \textbf{238} (1990) 92; hep-th/9704166.
\item Codello--Zanusso (CZ): arXiv:1203.2034.
\item Vassilevich: Phys.\ Rept.\ \textbf{388} (2003) 279
  [hep-th/0306138].
\item Iochum--Levy--Vassilevich (ILV): arXiv:1108.3749; 1201.6637.
\end{itemize}
\end{abstract}

\tableofcontents
\newpage

%%======================================================================
\section{Conventions}
\label{sec:conventions}
%%======================================================================

All conventions are fixed here and \textbf{must be used consistently}
throughout the NT-1 computation.  Where references differ, we note the
conversion.

\subsection{Spacetime and metric}

\begin{convention}[Metric signature]
\oursign\;
We use the \textbf{mostly-plus} Lorentzian signature:
\begin{equation}
  \eta_{\mu\nu} = \operatorname{diag}(-1,+1,+1,+1)\,.
  \label{eq:signature}
\end{equation}
For the heat kernel computation we Wick-rotate to \textbf{Euclidean
signature} $(+,+,+,+)$.  The Euclidean metric is denoted $g_{\mu\nu}$
(positive definite).
\end{convention}

\begin{remark}
Vassilevich~\cite{Vassilevich2003}, Codello--Zanusso~\cite{CZ2012},
and the BV series~\cite{BV1987,BV1990a} all work in Euclidean
signature.  No sign adjustment needed for formulas extracted from
these sources (they are already Euclidean).  The Wick rotation from
our Lorentzian convention is
$t \to -\ii\tau$, $g^{00}_L = -1 \to g^{00}_E = +1$.
\end{remark}

\begin{convention}[Spacetime dimension]
$d = 4$ for all final results.  Intermediate formulas are kept
general in $d$ where possible.
\end{convention}


\subsection{Curvature tensors}

\begin{convention}[Riemann tensor]
\oursign\;
\begin{equation}
  R^{\rho}{}_{\sigma\mu\nu}
  = \partial_\mu \Gamma^\rho_{\nu\sigma}
  - \partial_\nu \Gamma^\rho_{\mu\sigma}
  + \Gamma^\rho_{\mu\lambda}\Gamma^\lambda_{\nu\sigma}
  - \Gamma^\rho_{\nu\lambda}\Gamma^\lambda_{\mu\sigma}\,.
  \label{eq:riemann}
\end{equation}
This is the ``$+--+$'' convention in the MTW classification
(same as Wald, Vassilevich, Avramidi).
\end{convention}

\begin{convention}[Ricci tensor]
\oursign\;
Contraction on \textbf{first and third} indices:
\begin{equation}
  R_{\mu\nu} = R^{\rho}{}_{\mu\rho\nu}\,.
  \label{eq:ricci}
\end{equation}
This matches Vassilevich~\cite{Vassilevich2003} and
Codello--Zanusso~\cite{CZ2012} (who write
$R_{\mu\nu} \equiv R_{\alpha\mu}{}^{\alpha}{}_{\nu}$, which is
the same thing).
\end{convention}

\begin{convention}[Scalar curvature]
\begin{equation}
  R = g^{\mu\nu} R_{\mu\nu}\,.
\end{equation}
Positive for a round sphere in Euclidean signature.
\end{convention}

\begin{convention}[Weyl tensor]
In $d$ dimensions:
\begin{equation}
  C_{\mu\nu\rho\sigma}
  = R_{\mu\nu\rho\sigma}
  - \frac{2}{d-2}\bigl(
    g_{\mu[\rho}R_{\sigma]\nu} - g_{\nu[\rho}R_{\sigma]\mu}
  \bigr)
  + \frac{2}{(d-1)(d-2)}\,R\,g_{\mu[\rho}g_{\sigma]\nu}\,.
  \label{eq:weyl}
\end{equation}
Its square satisfies \fromlit{CZ~(A.3)}:
\begin{equation}
  \Weylsq = \Rsq - \frac{4}{d-2}\Ricsq + \frac{2}{(d-1)(d-2)}\,R^2\,.
  \label{eq:weyl-sq}
\end{equation}
In $d=4$: $\Weylsq = \Rsq - 2\Ricsq + \frac{1}{3}R^2$.
\end{convention}

\begin{convention}[Gauss--Bonnet / Euler density]
\begin{equation}
  E_4 = \Rsq - 4\Ricsq + R^2\,.
  \label{eq:euler}
\end{equation}
Topological in $d=4$: $\int E_4\sqrt{g}\,\dd^4x = 32\pi^2\chi(M)$.
\end{convention}


\subsection{Differential operators}

\begin{convention}[Covariant derivative and d'Alembertian]
\oursign\;
The covariant d'Alembertian (Laplace--Beltrami operator) on scalars:
\begin{equation}
  \Box \equiv g^{\mu\nu}\nabla_\mu\nabla_\nu\,.
  \label{eq:box}
\end{equation}
In Lorentzian signature this is the wave operator ($\Box\phi =
-\partial_t^2\phi + \nabla^2\phi$ in flat space).
In \textbf{Euclidean signature} this becomes:
\begin{equation}
  \Box_E = g^{\mu\nu}\nabla_\mu\nabla_\nu = \nabla^2
  \quad\text{(positive spectrum on compact manifold)}\,.
\end{equation}
\end{convention}

\begin{convention}[Generalised Laplacian]
\label{conv:laplacian}
\oursign\;
A Laplace-type operator on sections of a vector bundle:
\begin{equation}
  \boxed{
  \Delta = -(g^{\mu\nu}D_\mu D_\nu + E)
  = -D^2 - E
  }
  \label{eq:laplacian}
\end{equation}
where $D_\mu = \nabla_\mu + A_\mu$ is the full covariant derivative
(including gauge/spin connection) and $E$ is the endomorphism (potential
term).

This matches Vassilevich~\cite{Vassilevich2003} Eq.~(1.1) and
Codello--Zanusso~\cite{CZ2012} Eq.~(2).

\medskip\noindent
\textbf{Sign note:} CZ define $\Delta = -D^2 + \mathbf{U}$ with
$\mathbf{U}$ their endomorphism.  Comparing with our convention
$\Delta = -D^2 - E$, we identify:
\begin{equation}
  \mathbf{U}_{\text{CZ}} = -E\,.
  \label{eq:U-vs-E}
\end{equation}
This sign difference is critical when translating form factors.
\end{convention}

\begin{convention}[Bundle curvature]
The curvature of the connection $A_\mu$ on the vector bundle:
\begin{equation}
  \Omega_{\mu\nu} \equiv [D_\mu, D_\nu]
  = \partial_\mu A_\nu - \partial_\nu A_\mu + [A_\mu, A_\nu]\,.
  \label{eq:bundle-curv}
\end{equation}
This matches both CZ Eq.~(1) and Vassilevich.
\end{convention}

\begin{convention}[BV endomorphism $\hat{P}$]
\label{conv:Phat}
The BV papers use $\hat{P} \equiv E + \frac{R}{6}\Id$
(equivalently $\hat{P} = -\mathbf{U}_{\text{CZ}} + \frac{R}{6}\Id$)
\fromlit{Vassilevich Eq.~(4.3); CZ App.~A}.
In our conventions for the Dirac operator:
\begin{equation}
  E = -\frac{R}{4}\Id_4
  \quad\Longrightarrow\quad
  \hat{P} = -\frac{R}{4}\Id_4 + \frac{R}{6}\Id_4
  = -\frac{R}{12}\Id_4\,.
  \label{eq:Phat-dirac}
\end{equation}
The CZ relation is $\mathbf{U}_{\text{CZ}} = -\hat{P} + \frac{R}{6}\Id = -E$,
so $\hat{P} = -\mathbf{U}_{\text{CZ}} + \frac{R}{6}\Id$ \fromlit{CZ App.~A}.

\signcorrection{Convention~2.11 corrected 2026-03-09: the BV endomorphism
is $\hat{P} = E + R/6$ (not $E - R/6$ as stated in the original version).
The original error gave $\hat{P} = -5R/12$ instead of the correct
$\hat{P} = -R/12$ for Dirac.  This error was \emph{contained}:
the NT-1 derivation pipeline (Proposition~3.2, Theorem~4.1 of
\texttt{NT1\_form\_factors.tex}) works entirely through CZ form factors
and never uses $\hat{P}$ in the computational chain.  See the failed
derivation in Appendix~A of \texttt{NT1\_form\_factors.tex} for the
consequences of the wrong $\hat{P}$.}
\end{convention}

\begin{remark}[$\hat{P}$-independence of NT-1]
\label{rem:Phat-independence}
The NT-1 derivation of the nonlocal form factors $F_1$, $F_2$
is \emph{structurally independent} of the BV endomorphism~$\hat{P}$.
The computational pipeline proceeds as follows:
\begin{enumerate}[nosep]
  \item Extract operator data $(E, \Omega_{\mu\nu})$ from the
        Lichnerowicz formula (Convention~\ref{eq:lichnerowicz}).
  \item Map to the CZ universal form factors
        $f_{\mathrm{Ric}}, f_R, f_{R\mathbf{U}}, f_\mathbf{U}, f_\Omega$
        using the CZ notation
        $\mathbf{U}_{\mathrm{CZ}} = -E$ (Convention~2.10).
  \item Contract CZ form factors with Dirac traces
        $\tr(\Id)$, $\tr(E)$, $\tr(E^2)$,
        $\tr(\Omega_{\mu\nu}\Omega^{\mu\nu})$
        to obtain $h_C(x)$ and $h_R(x)$.
\end{enumerate}
At no point does $\hat{P}$ or its traces appear in this chain.
The BV form factors $f_1,\dots,f_5$ (which do involve $\hat{P}$)
are used only in \S3 of \texttt{NT1\_form\_factors.tex} as an
alternative cross-check pathway, which was found to be algebraically
more cumbersome and is not part of the main derivation.
\end{remark}


\subsection{Dirac operator specifics}

\begin{convention}[SCT Dirac operator]
On a closed Riemannian spin 4-manifold $(M,g)$:
\begin{equation}
  D = \ii\gamma^\mu\nabla^S_\mu\,,
  \label{eq:dirac}
\end{equation}
where $\nabla^S_\mu = \nabla_\mu + \frac{1}{4}\omega_\mu^{ab}\gamma_a\gamma_b$
is the spin connection covariant derivative.
\end{convention}

\begin{convention}[Lichnerowicz formula]
\begin{equation}
  D^2 = -g^{\mu\nu}\nabla^S_\mu\nabla^S_\nu + \frac{R}{4}\,.
  \label{eq:lichnerowicz}
\end{equation}
So $D^2$ is of Laplace type \eqref{eq:laplacian} with:
\begin{equation}
  \boxed{E = -\frac{R}{4}\Id_4\,,
  \qquad
  \Omega_{\mu\nu}^{(\text{spin})}
  = \frac{1}{4}R_{\mu\nu\rho\sigma}\gamma^\rho\gamma^\sigma}
  \label{eq:dirac-data}
\end{equation}
and $\mathbf{U}_{\text{CZ}} = -E = \frac{R}{4}\Id_4$.

\begin{remark}[\textbf{Sign of $E$ --- critical}]
\label{rem:E-sign}
The Lichnerowicz formula gives $D^2 = -g^{\mu\nu}\nabla\nabla + R/4$.
Comparing with the generalized Laplacian $\Delta = -(g^{\mu\nu}\nabla\nabla + E)$,
we identify $E = -R/4$ (\textbf{not} $+R/4$).
Cross-check: $E = -R/4$ gives $h_R(0) = 0$, confirming conformal invariance
of the massless Dirac field. The wrong sign $E = +R/4$ would give $h_R(0) \neq 0$.
\end{remark}
\end{convention}

\begin{convention}[Spin connection curvature trace identities]
\exactlit\;
\begin{align}
  \tr(\Omega_{\mu\nu}^{(\text{spin})}\Omega^{(\text{spin})\mu\nu})
  &= \frac{1}{16}\,R_{\mu\nu\rho\sigma}R_{\alpha\beta\gamma\delta}\,
    \tr(\gamma^\rho\gamma^\sigma\gamma^\gamma\gamma^\delta)\,
    g^{\mu\alpha}g^{\nu\beta}
  \notag\\
  &= -\frac{1}{2}\,\Rsq\,.
  \label{eq:OmOm-trace}
\end{align}
(The factor $N_{\text{spin}}=\tr\Id_4=4$ in $d=4$.)
\end{convention}


\subsection{Heat kernel and spectral action}

\begin{convention}[Heat kernel]
\begin{equation}
  K(s;x,y) = \langle x|e^{-s\Delta}|y\rangle\,,
  \qquad
  \Tr\, K(s) = \Tr\, e^{-s\Delta}\,.
  \label{eq:heat-kernel}
\end{equation}
Here $s > 0$ is the proper-time parameter.
The heat kernel satisfies $(\partial_s + \Delta_x)K(s;x,y)=0$
with $K(0;x,y) = \delta(x,y)$.
This matches CZ~\eqref{eq:laplacian}--\cite{CZ2012} Eqs.~(3)--(4).
\end{convention}

\begin{convention}[SCT spectral action and heat kernel]
\label{conv:spectral}
The spectral action is
\begin{equation}
  S_{\text{SCT}} = \Tr\!\bigl(f(D^2/\Lambda^2)\bigr)\,.
  \label{eq:spectral-action}
\end{equation}
Using the Laplace transform $f(u) = \int_0^\infty \tilde{f}(t)\,
e^{-tu}\dd t$:
\begin{equation}
  \Tr\!\bigl(f(D^2/\Lambda^2)\bigr)
  = \int_0^\infty \tilde{f}(t)\,
    \Tr\!\bigl(e^{-t D^2/\Lambda^2}\bigr)\dd t\,.
  \label{eq:spectral-heat}
\end{equation}
With $s = t/\Lambda^2$:
\begin{equation}
  \boxed{
  \Tr\!\bigl(f(D^2/\Lambda^2)\bigr)
  = \Lambda^2\int_0^\infty \tilde{f}(s\Lambda^2)\,
    \Tr\!\bigl(e^{-s D^2}\bigr)\dd s
  }
  \label{eq:spectral-heat-s}
\end{equation}
Note $D^2 = \Delta$ (our Laplace-type operator for the Dirac case).
\end{convention}

\begin{convention}[Relation to one-loop effective action]
The standard one-loop effective action for a field with operator $\Delta$
and mass $m$:
\begin{equation}
  \Gamma^{(1)} = \mp\frac{1}{2}\ln\det\Delta
  = \pm\frac{1}{2}\int_0^\infty \frac{\dd s}{s}\,
    e^{-sm^2}\,\Tr\,e^{-s\Delta}\,,
  \label{eq:eff-action}
\end{equation}
where $-$ for bosons, $+$ for fermions (Dirac).
Compare with \eqref{eq:spectral-heat-s}: the spectral action replaces
the $1/s$ weight with a smooth cutoff function.
\end{convention}


%%======================================================================
\newpage
\section{Avramidi Covariant Technique for Off-Diagonal Heat Kernel}
\label{sec:avramidi}
%%======================================================================

\subsection{The problem: local vs.\ nonlocal}

The Seeley--DeWitt expansion gives the heat trace as a series in
powers of proper time $s$:
\begin{equation}
  \Tr\,e^{-s\Delta} \sim \frac{1}{(4\pi s)^{d/2}}
  \sum_{k=0}^\infty s^k\, a_{2k}(\Delta)\,,
  \qquad s \to 0^+\,.
  \label{eq:SD-expansion}
\end{equation}
This is \textbf{local}: each $a_{2k}$ is a spacetime integral of local
curvature invariants.  It captures only the \textbf{coincidence limit}
$x=y$ of the heat kernel diagonal.

To obtain the \textbf{nonlocal form factors} $F_1(\Box/\Lambda^2)$,
$F_2(\Box/\Lambda^2)$ we need the heat kernel expanded to second
order in curvature but to \textbf{all orders in derivatives}
(i.e., all orders in $s$, not just the $s\to 0$ asymptotics).

\subsection{Avramidi's partial summation}
\label{subsec:avramidi-partial}

\fromlit{Avramidi, Phys.\ Lett.\ B \textbf{236} (1990) 443;
         Phys.\ Lett.\ B \textbf{238} (1990) 92;
         hep-th/9704166}

Avramidi's key insight: decompose the heat kernel expansion not by
powers of $s$ (local) but by \textbf{powers of curvature}, while
keeping all powers of $\nabla$ (i.e., all powers of $s\Box$).
This ``partial summation'' resums the Seeley--DeWitt series into
nonlocal form factors.

\begin{definition}[Avramidi form factor basis]
At second order in curvature, the nonlocal heat trace on a
boundaryless manifold takes the form
\fromlit{Avramidi (1990); CZ Eq.~(23)}:
\begin{equation}
  \boxed{
  \begin{aligned}
  \Tr\,e^{-s\Delta}
  &= \frac{1}{(4\pi s)^{d/2}}
  \int\dd^d x\,\sqrt{g}\,\tr\Bigl\{
  \Id + s\hat{P}
  + s^2\Bigl[
  \Id\,R_{\mu\nu}\, f_1(s\Box)\, R^{\mu\nu}\\
  &\quad+ \Id\,R\,f_2(s\Box)\,R
  + R\,f_3(s\Box)\,\hat{P}
  + \hat{P}\,f_4(s\Box)\,\hat{P}
  + \Omega_{\mu\nu}\,f_5(s\Box)\,\Omega^{\mu\nu}
  \Bigr]
  + \mathcal{O}(\mathcal{R}^3)
  \Bigr\}
  \end{aligned}
  }
  \label{eq:BV-heat-trace}
\end{equation}
where $\hat{P} = E - \frac{R}{6}\Id$ (Convention~\ref{conv:Phat}),
$\Box = -D^2$ (\textbf{without} the endomorphism),
and $f_i$ are the five BV form factors.
\end{definition}

\begin{remark}
This expansion is valid under the \textbf{covariant perturbation theory
condition} \fromlit{CZ~(24)}:
\begin{equation}
  \nabla\nabla\mathcal{R} \gg \mathcal{R}^2\,,
  \label{eq:cpt-condition}
\end{equation}
where $\mathcal{R}$ represents any curvature.  Physically: curvatures
are weak but may vary rapidly.  This is complementary to the
Seeley--DeWitt regime (strong/slowly varying fields).
\end{remark}


\subsection{The basic form factor (master function)}

\begin{definition}[Master form factor]
\exactlit\;
\fromlit{CZ Eq.~(36); Avramidi (1990); BV (1987)}
\begin{equation}
  \boxed{
  f(x) = \int_0^1 \dd\xi\; e^{-\xi(1-\xi)x}
  }
  \label{eq:master-f}
\end{equation}
where $x = s\Box$ (or $x = sp^2$ in momentum space).
\end{definition}

\paragraph{Properties:}
\begin{itemize}
\item \textbf{Normalisation:} $f(0) = 1$.
\item \textbf{Small-$x$ expansion} \fromlit{CZ below Eq.~(37)}:
  \begin{equation}
    f(x) = 1 - \frac{x}{6} + \frac{x^2}{60} - \frac{x^3}{840}
    + \mathcal{O}(x^4)\,.
    \label{eq:f-small}
  \end{equation}
  The general coefficient is
  $\frac{(-1)^n}{n!}B(n+1,n+1)$ where $B$ is the Beta function.
\item \textbf{Large-$x$ expansion} \fromlit{CZ below Eq.~(39)}:
  \begin{equation}
    f(x) = \frac{2}{x} + \frac{4}{x^2}
    + \mathcal{O}\!\left(\frac{1}{x^3}\right)\,,
    \qquad x\to+\infty\,.
    \label{eq:f-large}
  \end{equation}
\item \textbf{Saddle point:} For $x\to+\infty$,
  $f(x) \sim \sqrt{\pi/x}\,e^{-x/4}$
  (the saddle at $\xi=1/2$ gives the exponential tail).
\item \textbf{Useful identity} \fromlit{CZ Sec.~4.2}:
  \begin{equation}
    \frac{1}{2}\int_0^1\dd\xi\, e^{-\xi(1-\xi)x}(1-2\xi)^2
    = -\frac{1}{x}[f(x)-1]\,.
    \label{eq:f-identity}
  \end{equation}
\end{itemize}


%%======================================================================
\newpage
\section{Form Factor Construction: Heat Kernel}
\label{sec:form-factors-hk}
%%======================================================================

\subsection{The five BV form factors}

\fromlit{CZ Eqs.~(37)--(38); BV (1987) Eqs.~(3.44)--(3.48)}

\exactlit\;
In the ``Ricci basis'' $\{R_{\mu\nu}, R, \mathbf{U}, \Omega_{\mu\nu}\}$
used by CZ, the nonlocal heat trace at $\mathcal{O}(\mathcal{R}^2)$ is:
\begin{equation}
  \Tr\,e^{-s\Delta}
  = \frac{1}{(4\pi s)^{d/2}}
  \int\dd^d x\,\sqrt{g}\,\tr\Bigl\{
  \Id - s\mathbf{U} + s\Id\frac{R}{6}
  + s^2[\text{form factor terms}]
  + \mathcal{O}(\mathcal{R}^3)
  \Bigr\}
  \label{eq:CZ-heat-trace}
\end{equation}
with form factor terms:
\begin{equation}
  \Id\,R_{\mu\nu}\,f_{\text{Ric}}(s\Box)\,R^{\mu\nu}
  + \Id\,R\,f_R(s\Box)\,R
  + R\,f_{R\mathbf{U}}(s\Box)\,\mathbf{U}
  + \mathbf{U}\,f_\mathbf{U}(s\Box)\,\mathbf{U}
  + \Omega_{\mu\nu}\,f_\Omega(s\Box)\,\Omega^{\mu\nu}\,.
\end{equation}

The form factors in terms of the master function $f(x)$:
\begin{equation}
  \boxed{
  \begin{aligned}
  f_{\text{Ric}}(x)
    &= \frac{1}{6x} + \frac{1}{x^2}[f(x)-1]\\[4pt]
  f_R(x)
    &= \frac{1}{32}f(x) + \frac{1}{8x}f(x) - \frac{7}{48x}
       - \frac{1}{8x^2}[f(x)-1]\\[4pt]
  f_{R\mathbf{U}}(x)
    &= -\frac{1}{4}f(x) - \frac{1}{2x}[f(x)-1]\\[4pt]
  f_\mathbf{U}(x)
    &= \frac{1}{2}f(x)\\[4pt]
  f_\Omega(x)
    &= -\frac{1}{2x}[f(x)-1]
  \end{aligned}
  }
  \label{eq:CZ-form-factors}
\end{equation}

\fromlit{CZ Eq.~(37), confirmed to match BV (1987, 1990) via
basis change in CZ App.~A}


\subsection{Conversion to BV basis}

The BV papers use $\hat{P} = E - \frac{R}{6}\Id$ instead of $\mathbf{U}$,
with $\mathbf{U} = -E$ (Convention~\ref{conv:laplacian}).
The conversion \fromlit{CZ Eq.~(50)} is:
\begin{equation}
  \boxed{
  \begin{aligned}
  f_1 &= f_{\text{Ric}}\,,\\
  f_2 &= f_R + \frac{1}{6}f_{R\mathbf{U}} + \frac{1}{36}f_\mathbf{U}\,,\\
  f_3 &= -f_{R\mathbf{U}} - \frac{1}{3}f_\mathbf{U}\,,\\
  f_4 &= f_\mathbf{U}\,,\\
  f_5 &= f_\Omega\,.
  \end{aligned}
  }
  \label{eq:BV-conversion}
\end{equation}

Substituting \eqref{eq:CZ-form-factors} into \eqref{eq:BV-conversion}
gives the explicit BV form factors.  This can be verified against
BV~(1987) Table~1.


\subsection{Conversion to Weyl basis}

\fromlit{CZ Eqs.~(41)--(42); Barvinsky--Mirzabekian--Zhytnikov gr-qc/9510037}

For the Weyl basis $\{C_{\mu\nu\rho\sigma}, R, \hat{P}, \Omega_{\mu\nu}\}$:
\begin{equation}
  \boxed{
  \begin{aligned}
  f_C(x) &= \frac{d-2}{4(d-3)}\,f_{\text{Ric}}(x)
  \qquad\text{(in $d=4$: $f_C = \tfrac{1}{2}f_{\text{Ric}}$)}\\[4pt]
  f_{R,\text{bis}}(x) &= \frac{d}{4(d-1)}\,f_{\text{Ric}}(x) + f_R(x)
  \qquad\text{(in $d=4$: $f_{R,\text{bis}} = \tfrac{1}{3}f_{\text{Ric}} + f_R$)}
  \end{aligned}
  }
  \label{eq:Weyl-conversion}
\end{equation}

The heat trace in the Weyl basis \fromlit{CZ Eq.~(40)}:
\begin{equation}
  \Tr\,e^{-s\Delta}
  = \frac{1}{(4\pi s)^{d/2}}
  \int\dd^d x\,\sqrt{g}\,\tr\Bigl\{
  \Id - s\mathbf{U} + s\frac{R}{6}\Id
  + s^2\Bigl[
  \Id\,C_{\mu\nu\rho\sigma}\,f_C\,C^{\mu\nu\rho\sigma}
  + \Id\,R\,f_{R,\text{bis}}\,R
  + \cdots
  \Bigr]\Bigr\}
  \label{eq:Weyl-heat-trace}
\end{equation}
where $\cdots$ contains $\hat{P}$ and $\Omega$ terms unchanged.

\begin{remark}[Universality]
In the Ricci basis, the form factors are \textbf{universal} (independent
of $d$ apart from the overall $(4\pi s)^{-d/2}$ prefactor).
The Weyl basis spoils universality by introducing explicit $d$-dependent
prefactors.
\end{remark}


\subsection{Short-time (local) limits}

\fromlit{CZ Eq.~(38)}

Setting $x=0$ in \eqref{eq:CZ-form-factors}:
\begin{equation}
  \begin{aligned}
  f_{\text{Ric}}(0) &= \frac{1}{60}\,,\\
  f_R(0) &= \frac{1}{120}\,,\\
  f_{R\mathbf{U}}(0) &= -\frac{1}{6}\,,\\
  f_\mathbf{U}(0) &= \frac{1}{2}\,,\\
  f_\Omega(0) &= \frac{1}{12}\,.
  \end{aligned}
  \label{eq:local-limits}
\end{equation}
These values reproduce the Seeley--DeWitt coefficient $a_4$ when
inserted into the heat trace.

\begin{remark}[Cross-check with $a_4$]
The local limit of \eqref{eq:CZ-heat-trace} must reproduce the
standard $a_4$ coefficient \fromlit{Vassilevich Eq.~(4.1)}:
\begin{equation}
  a_4 = \frac{1}{360}\int\dd^d x\,\sqrt{g}\,\tr\Bigl[
  60\,RE + 180\,E^2 + 60\,\Box E
  + (5R^2 - 2\Ricsq + 2\Rsq)\Id
  + 30\,\Omega_{\mu\nu}\Omega^{\mu\nu}
  \Bigr]\,.
  \label{eq:a4-vassilevich}
\end{equation}
This is identical to \eqref{eq:a4-V} in Appendix~\ref{sec:V}.
Note: Vassilevich uses $\Delta = -(g^{\mu\nu}\nabla_\mu\nabla_\nu + E)$
which is the \textbf{same} as our convention \eqref{eq:laplacian},
so his $E$ is the same as our $E$.
The CZ convention has $\mathbf{U} = -E$, so the signs of the $E$ terms
need careful tracking when converting.  The signs $+60\Box E$ and
$+30\,\Omega\Omega$ in Vassilevich's convention yield $-15\Rsq$
for the Dirac operator since
$\tr(\Omega^{\text{spin}}\Omega^{\text{spin}}) = -\frac{1}{2}\Rsq$.
\end{remark}


\subsection{Late-time (IR) limits}

\fromlit{CZ Eq.~(39)}

\begin{equation}
  \begin{aligned}
  f_{\text{Ric}}(x) &= \frac{1}{6x} - \frac{1}{x^2}
    + \mathcal{O}(x^{-3})\,,\\
  f_R(x) &= -\frac{1}{12x} + \frac{1}{2x^2}
    + \mathcal{O}(x^{-3})\,,\\
  f_{R\mathbf{U}}(x) &= -\frac{2}{x^2} + \mathcal{O}(x^{-3})\,,\\
  f_\mathbf{U}(x) &= \frac{1}{x} + \frac{2}{x^2}
    + \mathcal{O}(x^{-3})\,,\\
  f_\Omega(x) &= \frac{1}{2x} - \frac{1}{2x^2}
    + \mathcal{O}(x^{-3})\,.
  \end{aligned}
  \label{eq:late-time}
\end{equation}

These control the UV behavior of the effective action (large $\Box$).


%%======================================================================
\newpage
\section{BV Nonlocal Effective Action}
\label{sec:BV-action}
%%======================================================================

\subsection{From heat kernel to effective action}

\fromlit{BV (1987) Sec.~4; CZ Sec.~6}

The one-loop effective action for a field with Laplace-type operator
$\Delta$ and mass $m$:
\begin{equation}
  \Gamma^{(1)} = \frac{1}{2}\int_0^\infty\frac{\dd s}{s}\,
  e^{-sm^2}\,\Tr\,e^{-s\Delta}
  \label{eq:eff-action-heat}
\end{equation}
(sign for bosons; the mass term provides IR regularisation).

Substituting the nonlocal heat trace \eqref{eq:CZ-heat-trace} and
performing the $s$-integration converts the heat kernel form factors
$f_i(s\Box)$ into effective action form factors $\mathcal{F}_i(\Box/m^2)$:
\begin{equation}
  \mathcal{F}_i(\Box/m^2)
  = \int_0^\infty \frac{\dd s}{s}\,e^{-sm^2}\,(4\pi s)^{-d/2}\,
  s^2\, f_i(s\Box)\,.
  \label{eq:eff-action-ff}
\end{equation}

\begin{remark}
The factor $s^2$ comes from the fact that the form factor terms in
\eqref{eq:CZ-heat-trace} carry an explicit $s^2$ prefactor at
$\mathcal{O}(\mathcal{R}^2)$.  The $1/s$ from
\eqref{eq:eff-action-heat} combines with this to give $s$.
\end{remark}


\subsection{Canonical form of the nonlocal effective action}

\fromlit{BV (1990) Eq.~(1.1); Barvinsky--Gusev--Vilkovisky--Zhytnikov
hep-th/9404187}

At $\mathcal{O}(\mathcal{R}^2)$, the one-loop effective action takes
the canonical form in the Weyl basis:
\begin{equation}
  \boxed{
  \Gamma^{(1)}_{\text{NL}}
  = \frac{1}{2(4\pi)^2}\int\dd^4 x\,\sqrt{g}\,\Bigl[
  C_{\mu\nu\rho\sigma}\,\mathcal{F}_W(\Box/\mu^2)\,C^{\mu\nu\rho\sigma}
  + R\,\mathcal{F}_R(\Box/\mu^2)\,R
  \Bigr]
  }
  \label{eq:BV-eff-action}
\end{equation}
(dropping the topological Gauss--Bonnet term in $d=4$, and the $\hat{P}$
and $\Omega$ terms which are absorbed into the Weyl/Ricci structures
for specific field content).


\subsection{Massless limit: logarithmic form factors}

\fromlit{BV (1990) Sec.~6; Barvinsky--Wachowski 2306.03780 Sec.~2}

In the massless limit, the effective action form factors develop
logarithmic behaviour:
\begin{equation}
  \mathcal{F}_i(\Box/\mu^2)
  \to \beta_i\,\ln(-\Box/\mu^2)
  + \text{finite local terms}\,.
  \label{eq:log-ff}
\end{equation}
The coefficients $\beta_i$ are scheme-independent and related to the
conformal anomaly.  This gives:
\begin{equation}
  \boxed{
  \Gamma^{(1)}_{\text{NL,massless}}
  = \frac{1}{32\pi^2}\int\dd^4 x\,\sqrt{g}\,\Bigl[
  \beta_W\, C_{\mu\nu\rho\sigma}\,
    \ln(-\Box/\mu^2)\, C^{\mu\nu\rho\sigma}
  + \beta_R\, R\,\ln(-\Box/\mu^2)\, R
  \Bigr]
  }
  \label{eq:log-action}
\end{equation}

\begin{table}[h]
\centering
\renewcommand{\arraystretch}{1.3}
\begin{tabular}{@{}lcc@{}}
\toprule
\textbf{Field type} & $\boldsymbol{\beta_W}$ & $\boldsymbol{\beta_R}$ \\
\midrule
Conformal scalar  & $1/120$  & $0$    \\
Dirac fermion     & $1/20$   & $0$    \\
Gauge vector      & $1/10$   & $0$    \\
Graviton          & $7/20$   & $0$    \\
\bottomrule
\end{tabular}
\caption{$\beta$-coefficients for conformally invariant fields.
\fromlit{BV (1990); Barvinsky--Vilkovisky--Zhytnikov hep-th/9404187}}
\label{tab:betas}
\end{table}

\begin{remark}
For all conformally invariant fields, $\beta_R = 0$.
This is the deep connection between conformal symmetry
and the structure of the nonlocal effective action.
\end{remark}


\subsection{The key formula: spectral action to form factors}

\gap\; \textbf{This is the central gap that NT-1 must bridge.}

The SCT spectral action $\Tr(f(D^2/\Lambda^2))$ gives a different
weight than the one-loop effective action $\frac{1}{2}\int\frac{\dd s}{s}
e^{-sm^2}\Tr\,e^{-sD^2}$.  Specifically:
\begin{equation}
  \Tr\!\bigl(f(D^2/\Lambda^2)\bigr)
  = \int_0^\infty \tilde{f}(t)\,\Tr\!\bigl(e^{-tD^2/\Lambda^2}\bigr)\dd t\,.
  \label{eq:spectral-vs-eff}
\end{equation}

Using \eqref{eq:BV-heat-trace}, the $\mathcal{O}(\mathcal{R}^2)$
contribution to the spectral action is:
\begin{equation}
  \Tr\!\bigl(f(D^2/\Lambda^2)\bigr)\Big|_{\mathcal{R}^2}
  = \int_0^\infty \tilde{f}(t)\,\frac{1}{(4\pi t/\Lambda^2)^2}
  \int\dd^4x\sqrt{g}\,\tr\bigl[
  \cdots\, s^2 f_i(s\Box)\,\cdots
  \bigr]_{s=t/\Lambda^2}\dd t
  \label{eq:spectral-R2}
\end{equation}
where the form factors $f_i$ are those of \eqref{eq:CZ-form-factors}.

The SCT form factors $F_1(\Box/\Lambda^2)$ and $F_2(\Box/\Lambda^2)$
are then obtained by performing the $t$-integration against the spectral
weight $\tilde{f}(t)$:
\begin{equation}
  \boxed{
  F_i^{\text{SCT}}(\Box/\Lambda^2)
  = \frac{\Lambda^4}{(4\pi)^2}\int_0^\infty \dd t\,
    \tilde{f}(t)\, (t/\Lambda^2)^{2-d/2}\,
    f_i(t\Box/\Lambda^2)
  }
  \label{eq:SCT-form-factor}
\end{equation}
where $f_i$ are the heat kernel form factors from
\eqref{eq:CZ-form-factors} or \eqref{eq:Weyl-conversion}.

\begin{remark}[Notation for the target]
In the Weyl basis, the final result is written as:
\begin{equation}
  S_{\text{SCT}}\Big|_{\mathcal{R}^2}
  = \int\dd^4x\,\sqrt{g}\,\Bigl[
  F_1(\Box/\Lambda^2)\,\Weylsq
  + F_2(\Box/\Lambda^2)\,R^2
  \Bigr]\,,
  \label{eq:target}
\end{equation}
where $F_1 \propto f_C$ integrated against $\tilde{f}$,
and $F_2 \propto f_{R,\text{bis}}$ integrated against $\tilde{f}$
(plus contributions from $\hat{P}$ and $\Omega$ terms traced for
the Dirac operator).
\end{remark}


%%======================================================================
\newpage
\section{Computational Recipe}
\label{sec:recipe}
%%======================================================================

The step-by-step procedure for computing $F_1(\Box/\Lambda^2)$ and
$F_2(\Box/\Lambda^2)$ from the SCT spectral action.

\subsection{Step 1: Identify the operator}

The SCT Dirac operator gives $D^2 = \Delta$ with:
\begin{itemize}
\item Endomorphism: $E = -\frac{R}{4}\Id_4$
  $\;\Longrightarrow\;$ $\mathbf{U}_{\text{CZ}} = -E = \frac{R}{4}\Id_4$
\item Bundle curvature:
  $\Omega_{\mu\nu} = \frac{1}{4}R_{\mu\nu\rho\sigma}\gamma^\rho\gamma^\sigma$
\item BV endomorphism:
  $\hat{P} = E - \frac{R}{6}\Id_4 = -\frac{5R}{12}\Id_4$
\item Spin dimension: $N = \tr\Id_4 = 4$
\end{itemize}


\subsection{Step 2: Write the nonlocal heat trace for $D^2$}

Using \eqref{eq:BV-heat-trace} with the Dirac data above.
The CZ form factors $f_1, \ldots, f_5$ act on the CZ endomorphism
$\mathbf{U}_{\text{CZ}} = -E = \frac{R}{4}\Id_4$, \textbf{not} on
$\hat{P}$.  Writing the heat trace in CZ notation:
\begin{equation}
  \begin{aligned}
  \Tr\,e^{-sD^2}
  &= \frac{1}{(4\pi s)^2}
  \int\dd^4 x\,\sqrt{g}\,\Bigl\{
  4 - s\,\frac{5R}{3}\\
  &\quad+ s^2\Bigl[
  4\,R_{\mu\nu}\,f_1(s\Box)\,R^{\mu\nu}
  + 4\,R\,f_2(s\Box)\,R\\
  &\quad+ R\,f_3(s\Box)\,\tr(\mathbf{U}_{\text{CZ}})
  + f_4(s\Box)\,\tr(\mathbf{U}_{\text{CZ}}^2)
  + f_5(s\Box)\,\tr(\Omega\Omega)
  \Bigr]
  + \mathcal{O}(\mathcal{R}^3)
  \Bigr\}
  \end{aligned}
  \label{eq:heat-trace-dirac}
\end{equation}

After taking the trace over spinor indices (with $E = -R/4$,
$\mathbf{U}_{\text{CZ}} = R/4$, $\hat{P} = -5R/12$):
\begin{itemize}
\item $\tr(\mathbf{U}_{\text{CZ}}) = 4\cdot R/4 = R$
\item $\tr(\mathbf{U}_{\text{CZ}}^2) = 4\cdot R^2/16 = R^2/4$
\item $\tr(\hat{P}) = 4\cdot(-5R/12) = -5R/3$
\item $\tr(\hat{P}^2) = 4\cdot 25R^2/144 = 25R^2/36$
\item $\tr(\Omega_{\mu\nu}\Omega^{\mu\nu}) = -\frac{1}{2}\Rsq$
\end{itemize}

So the $\mathcal{O}(\mathcal{R}^2)$ part becomes:
\begin{equation}
  \begin{aligned}
  \Tr\,e^{-sD^2}\Big|_{\mathcal{R}^2}
  &= \frac{s^2}{(4\pi s)^2}\int\dd^4x\,\sqrt{g}\,\Bigl[
  4\,R_{\mu\nu}\,f_1(s\Box)\,R^{\mu\nu}\\
  &\quad+ \Bigl(4f_2 + f_3 + \frac{1}{4}f_4\Bigr)\,
    R\,(s\Box)\,R
  - \frac{1}{2}\,\Rsq\,f_5(s\Box)
  \Bigr]
  \end{aligned}
  \label{eq:heat-trace-dirac-R2}
\end{equation}


\subsection{Step 3: Convert to Weyl basis}

Using the Gauss--Bonnet identity at $\mathcal{O}(\mathcal{R}^2)$
\fromlit{CZ App.~A; Avramidi (2000)}:
\begin{equation}
  \int\dd^4x\,\sqrt{g}\, R_{\mu\nu\rho\sigma}\,h(\Box)\,R^{\mu\nu\rho\sigma}
  = \int\dd^4x\,\sqrt{g}\,\bigl[
  4\,R_{\mu\nu}\,h(\Box)\,R^{\mu\nu} - R\,h(\Box)\,R
  \bigr] + \mathcal{O}(\mathcal{R}^3)\,.
  \label{eq:GB-nonlocal}
\end{equation}
This means the $\Rsq$ term from $\tr(\Omega\Omega)$ can be absorbed
into $\Ricsq$ and $R^2$ terms.

Then convert $\Ricsq \to \Weylsq + \frac{2}{d-2}\Ricsq_{\text{tracefree}} + \ldots$
using \eqref{eq:weyl-sq}.

The final result should take the form:
\begin{equation}
  \Tr\,e^{-sD^2}\Big|_{\mathcal{R}^2}
  = \frac{s^2}{(4\pi s)^2}\int\dd^4x\,\sqrt{g}\,\Bigl[
  h_C(s\Box)\,\Weylsq + h_R(s\Box)\,R^2
  \Bigr]
  \label{eq:heat-trace-weyl}
\end{equation}
where $h_C$ and $h_R$ are combinations of the $f_i$ with the
appropriate Dirac trace coefficients.


\subsection{Step 4: Integrate against spectral weight}

The SCT form factors are:
\begin{equation}
  \boxed{
  \begin{aligned}
  F_1(\Box/\Lambda^2) &= \frac{\Lambda^4}{(4\pi)^2}
    \int_0^\infty\dd t\,\tilde{f}(t)\,
    h_C(t\Box/\Lambda^2)\\[6pt]
  F_2(\Box/\Lambda^2) &= \frac{\Lambda^4}{(4\pi)^2}
    \int_0^\infty\dd t\,\tilde{f}(t)\,
    h_R(t\Box/\Lambda^2)
  \end{aligned}
  }
  \label{eq:F1F2-integral}
\end{equation}
where the $t$-integration converts heat kernel form factors to spectral
action form factors.

\begin{remark}[Difference from one-loop effective action]
In the standard one-loop computation, the weight is
$\frac{1}{s}e^{-sm^2}$ (gives logarithmic form factors in massless limit).
In SCT, the weight is $\tilde{f}(t)$ (the inverse Laplace transform
of the spectral function $f$), which provides UV regularisation at
scale $\Lambda$.  The SCT form factors are therefore
\textbf{manifestly finite} --- no renormalisation needed.
\end{remark}


\subsection{Step 5: Check limiting cases}

\begin{enumerate}
\item \textbf{Local limit} ($\Box \to 0$):
  $F_1(0)$ and $F_2(0)$ should reproduce the $c_1$ and $c_2$ coefficients
  from the Seeley--DeWitt computation (Derivation~1 in
  \texttt{SCT\_derivations.tex}).

\item \textbf{Flat space limit} ($R \to 0$):
  Both $F_1$ and $F_2$ should vanish trivially.

\item \textbf{Large-$\Box$ limit} ($\Box/\Lambda^2 \gg 1$):
  Form factors should decay (UV behaviour controlled by the spectral
  function $f$).

\item \textbf{Comparison with $\beta$-coefficients}:
  In the limit where $f$ approaches a sharp cutoff, the form factors
  should reproduce the logarithmic running with the correct
  $\beta_W = 1/20$ for the Dirac field.
\end{enumerate}


%%======================================================================
\newpage
\section{Key Formulas from Literature}
\label{sec:key-formulas}
%%======================================================================

\subsection{Codello--Zanusso diagrammatic computation}

\fromlit{arXiv:1203.2034, Secs.~3--5}

The CZ paper provides an independent re-derivation of the BV form
factors using a diagrammatic expansion of the heat kernel.  Key results:

\begin{enumerate}
\item \textbf{Perturbative expansion of heat kernel}
\fromlit{CZ Eq.~(16)}:
\begin{equation}
  K_{xy}^s = K_{0,xy}^s
  - s\int_0^1\dd t\int_z K_{0,xz}^{s(1-t)}V_z K_{0,zy}^{st}
  + \mathcal{O}(V^2)
  \label{eq:CZ-perturbative}
\end{equation}
where $K_{0,xy}^s = (4\pi s)^{-d/2}e^{-(x-y)^2/4s}$ is the flat
space heat kernel, and $V$ contains all interactions
(metric fluctuation $h_{\mu\nu}$, connection $A_\mu$, and
endomorphism $\mathbf{U}$).

\item \textbf{General $n$-th order term} \fromlit{CZ Eq.~(17)}:
\begin{equation}
  (-s)^n \int_0^1\dd t_1\cdots\int_0^{t_{n-1}}\dd t_n
  \int_{z_1\cdots z_n}
  K_{0,xz_1}^{s(1-t_1)}V_{z_1}\cdots V_{z_n}K_{0,z_ny}^{st_n}
  \label{eq:CZ-nth-order}
\end{equation}

\item \textbf{Master function from momentum integral}
\fromlit{CZ Eq.~(36)}:

The basic form factor emerges from the sunset diagram:
\begin{equation}
  f(sp^2) = \int_0^1\dd\xi\,e^{-\xi(1-\xi)sp^2}
  \label{eq:CZ-master-from-diagram}
\end{equation}
where $\xi = t_1 - t_2$ is a reparametrisation of the proper-time
integration variables.

\item \textbf{All five form factors confirmed} \fromlit{CZ Sec.~5}:
The CZ computation reproduces the BV form factors in
\eqref{eq:CZ-form-factors} by explicit diagrammatic evaluation
of 1-point and 2-point functions of the heat trace.
\end{enumerate}


\subsection{Avramidi generating function method}

\fromlit{Avramidi, hep-th/9704166 Secs.~2--4;
         Phys.\ Lett.\ B \textbf{238} (1990) 92}

\begin{enumerate}
\item \textbf{Covariantly constant curvature:}
Avramidi obtains closed-form heat kernel for covariantly constant
backgrounds ($\nabla_\mu R_{\nu\rho\sigma\tau} = 0$,
$\nabla_\mu\Omega_{\nu\rho} = 0$, $\nabla_\mu E = 0$)
using a Lie-algebraic method.  The heat kernel diagonal is expressed
as an average over a Lie group with Gaussian measure
\fromlit{gr-qc/9408028, hep-th/9602169}.

\item \textbf{Partial summation:}
By expanding around the covariantly constant case and keeping
all orders in $s\Box$ while truncating at second order in curvature,
Avramidi obtains the same form factor structure as BV, confirming
the five universal form factors $f_1,\ldots,f_5$.

\item \textbf{Off-diagonal heat kernel:}
\fromlit{CZ App.~B; Avramidi (2000) Ch.~9}

The off-diagonal heat kernel (before taking the coincidence limit)
has additional form factors $g_{\mathbf{U}}(x)$ and $g_R(x)$:
\begin{equation}
  K^s(x,x) = \frac{\sqrt{g}}{(4\pi s)^{d/2}}\bigl\{
  \Id + s\,g_{\mathbf{U}}(s\Box)\,\mathbf{U}
  + s\,\Id\,g_R(s\Box)\,R
  + \mathcal{O}(\mathcal{R}^2)\bigr\}
  \label{eq:offdiag-hk}
\end{equation}
with \fromlit{CZ App.~B}:
\begin{equation}
  g_{\mathbf{U}}(x) = -f(x)\,,
  \qquad
  g_R(x) = -\frac{1}{2x} + \frac{1}{4x}(x+2)f(x)\,.
  \label{eq:offdiag-ff}
\end{equation}
These satisfy $g_{\mathbf{U}}(0) = -1$ and $g_R(0) = 1/6$.
\end{enumerate}


\subsection{Vassilevich: Seeley--DeWitt coefficients}

\fromlit{Vassilevich, hep-th/0306138, Sec.~4}

The standard Seeley--DeWitt coefficients for a Laplace-type operator
$\Delta = -(g^{\mu\nu}\nabla_\mu\nabla_\nu + E)$ on a closed manifold:

\begin{align}
  a_0 &= \frac{1}{(4\pi)^{d/2}}\int\dd^d x\,\sqrt{g}\,\tr\Id
  \label{eq:a0-V}\\[6pt]
  a_2 &= \frac{1}{(4\pi)^{d/2}}\frac{1}{6}
  \int\dd^d x\,\sqrt{g}\,\tr\bigl(6E + R\Id\bigr)
  \label{eq:a2-V}\\[6pt]
  a_4 &= \frac{1}{(4\pi)^{d/2}}\frac{1}{360}
  \int\dd^d x\,\sqrt{g}\,\tr\bigl[
  60\,RE + 180\,E^2 + 60\,\Box E
  \notag\\
  &\qquad + (5R^2 - 2\Ricsq + 2\Rsq)\Id
  + 30\,\Omega_{\mu\nu}\Omega^{\mu\nu}
  \bigr]
  \label{eq:a4-V}
\end{align}

\begin{remark}[Sign in $a_4$]
\needscheck\;
Vassilevich Eq.~(4.1) writes $a_4$ with $+180E^2$ and $+30\,\Omega\Omega$.
Some references (e.g., Gilkey) write $-180E^2$ and $-30\,\Omega\Omega$
due to a different sign convention for $E$.  Our
$E_{\text{Vassilevich}} = E_{\text{ours}}$ (both use $\Delta = -\nabla^2 - E$),
so \eqref{eq:a4-V} applies directly.
However, the sign of the $\Omega\Omega$ term depends on the convention
for the bundle curvature: if $\Omega_{\mu\nu} = [D_\mu,D_\nu]$
(our convention, same as CZ and Vassilevich), then
$\tr(\Omega\Omega)$ for the spin connection is \textbf{negative}:
$\tr(\Omega^{\text{spin}}\Omega^{\text{spin}}) = -\frac{1}{2}\Rsq$.
So $+30\times(-\frac{1}{2}\Rsq) = -15\Rsq$ enters $a_4$.
\end{remark}


\subsection{Iochum--Levy--Vassilevich: Spectral action beyond weak field}

\fromlit{arXiv:1108.3749; arXiv:1201.6637}

ILV compute the spectral action $\Tr\,f(D^2/\Lambda^2)$ covariantly
to second order in the gauge connection, going beyond the weak-field
(Seeley--DeWitt) approximation.  Key results:

\begin{enumerate}
\item They confirm that the covariant perturbation theory (BV technique)
applies to the spectral action: the form factors are the same as
those in the heat kernel, just integrated against a different weight.

\item For a commutative spectral triple on a torus, they compute
the spectral action to second order in $A_\mu$ and find
\fromlit{1108.3749 abstract}:
in the UV regime $p\to\infty$, the action decays as $1/p^4$ in any
even dimension.

\item The convergence of the nonlocal expansion is controlled by the
smoothness of $f$: for Schwartz-class $f$, the expansion converges
for all $\Box$.
\end{enumerate}

\gap\; \textbf{The explicit form of the spectral action form factors
(as opposed to the heat kernel form factors) for a general spectral
function $f$ and the full Dirac operator on a curved manifold has
not been computed in the literature.}  This is precisely the gap
NT-1 aims to fill.


\subsection{Barvinsky--Vilkovisky: parametric integral representation}

\fromlit{BV (1987) Eq.~(3.44) ff.; BV (1990) Eqs.~(3.1)--(3.4)}

The BV form factors have the parametric integral representation:
\begin{equation}
  f_i(x) = \int_0^1\dd\xi\,\Phi_i(\xi)\,e^{-\xi(1-\xi)x}
  \label{eq:BV-parametric}
\end{equation}
where the $\Phi_i(\xi)$ are polynomial weight functions.  For the
five form factors at $\mathcal{O}(\mathcal{R}^2)$:
\begin{equation}
  \boxed{
  \begin{aligned}
  \Phi_1(\xi) &= \frac{1}{6} + (1-6\xi+6\xi^2)\frac{1}{[\xi(1-\xi)]^2}
    \cdot\text{(subtracted)}\\
  \Phi_4(\xi) &= \frac{1}{2}\\
  \Phi_5(\xi) &= \frac{1}{2}\cdot\frac{(1-2\xi)^2}{[\text{from identity \eqref{eq:f-identity}}]}
  \end{aligned}
  }
  \label{eq:BV-weights}
\end{equation}

\gap\; \textbf{The exact polynomial weight functions $\Phi_1$, $\Phi_2$,
$\Phi_3$ from BV~(1987) are needed.}  The forms given above are
incomplete.  The original BV papers (Nucl.\ Phys.\ B \textbf{282}
and \textbf{333}) should be consulted for the precise expressions.
The arXiv paper 0911.1168 (BV~IV) contains the complete results
through third order in curvature.

\begin{remark}
The parametric representation \eqref{eq:BV-parametric} is particularly
useful for the NT-1 computation because it allows the $s$-integration
(Step~4 in Sec.~\ref{sec:recipe}) to be performed by interchanging
the order of integration:
\begin{equation}
  F_i^{\text{SCT}} \propto \int_0^1\dd\xi\,\Phi_i(\xi)
  \int_0^\infty\dd t\,\tilde{f}(t)\,e^{-\xi(1-\xi)t\Box/\Lambda^2}\,.
  \label{eq:interchange}
\end{equation}
The inner integral is then a function of $\xi(1-\xi)\Box/\Lambda^2$
determined by the spectral function $f$.
\end{remark}


%%======================================================================
\newpage
\section{Known Results for Cross-Checks}
\label{sec:cross-checks}
%%======================================================================

\subsection{Seeley--DeWitt limit}

In the local limit $\Box/\Lambda^2 \to 0$, the form factors $F_1$ and
$F_2$ should reduce to the curvature-squared coefficients obtained
from the $a_4$ Seeley--DeWitt coefficient.

From \texttt{SCT\_derivations.tex} Eq.~(1.22):
\begin{equation}
  a_4(D^2) = \frac{1}{16\pi^2\cdot 360}\int_M\bigl[
  12\,\Box R + 5R^2 - 8\Ricsq - 7\Rsq
  \bigr]\sqrt{g}\,\dd^4x
  \label{eq:a4-check}
\end{equation}

In the Weyl basis: $5R^2 - 8\Ricsq - 7\Rsq = 11\,E_4 - 18\,C^2$.

So the local limit of the spectral action at $\mathcal{O}(\mathcal{R}^2)$ gives:
\begin{equation}
  F_1(0) = -\frac{18\,f_4}{5760\pi^2} = -\frac{f_4}{320\pi^2}\,,
  \qquad
  F_2(0) = 0\,,
  \label{eq:local-check}
\end{equation}
where $f_4 = f(0)$ is the zeroth moment of the spectral function.
The vanishing $F_2(0) = 0$ is a consequence of conformal invariance of the
massless Dirac field: $\beta_R = 0$ (verified independently in
NT1\_form\_factors, Sec.~5).

\begin{remark}
The $E_4$ (topological) part does not contribute to $F_1$ or $F_2$
in the equations of motion, but enters the action.  When checking
the local limit, one should compare with the full $a_4$ including
$E_4$.  The $\Box R$ term integrates to zero on a closed manifold.
\end{remark}


\subsection{Conformal flatness limit}

On a conformally flat manifold ($C_{\mu\nu\rho\sigma}=0$), only
the $F_2\,R^2$ term survives.  This should be consistent with the
known spectral action on spaces like $S^4$.


\subsection{Weak-field flat-space limit}

Expanding around flat space $g_{\mu\nu} = \delta_{\mu\nu} + h_{\mu\nu}$,
the form factors in momentum space should reproduce the 2-point
functions computed by CZ:
\begin{equation}
  \left.\frac{\delta^2\Tr K^s}{\delta h_{p;\mu\nu}\delta h_{-p;\alpha\beta}}
  \right|_{\Phi=0}
  = \frac{1}{(4\pi s)^{d/2}}\bigl\{
  \text{known tensor structure from CZ~(32)}
  \bigr\}
  \label{eq:flat-check}
\end{equation}


\subsection{Two-dimensional limit}

In $d=2$, the heat kernel expansion terminates at second order in
curvature \fromlit{CZ Sec.~5.1; BV (1990)}.
The combined form factor is \fromlit{CZ Eq.~(43)}:
\begin{equation}
  f_{R,2d}(x) = \frac{1}{32}f(x)
  + \frac{1}{16x}[2f(x)-1]
  + \frac{3}{8x^2}[f(x)-1]\,.
  \label{eq:2d-ff}
\end{equation}
This has the late-time expansion $f_{R,2d}(x) = 2/x^3 + \mathcal{O}(x^{-4})$,
with both $x^{-1}$ and $x^{-2}$ coefficients vanishing --- a nontrivial
check.


\subsection{Comparison with logarithmic running}

For a sharp cutoff spectral function (step function), the spectral
action form factors should reduce in the UV to:
\begin{equation}
  F_i(\Box/\Lambda^2)
  \xrightarrow{\Box/\Lambda^2 \to 0}
  \beta_i\,\ln(\Lambda^2/\Box) + \text{const}\,,
  \label{eq:log-check}
\end{equation}
matching the standard running with the correct $\beta$-coefficients
from Table~\ref{tab:betas}.

For the Dirac field: $\beta_W = 1/20$.  This provides a critical
cross-check of the final result.


\subsection{UV decay}

\fromlit{ILV 1108.3749}

For a smooth spectral function $f$ (Schwartz class), the spectral
action should decay in the UV ($p \to \infty$) as $1/p^4$ in $d=4$:
\begin{equation}
  F_i(p^2/\Lambda^2) \sim \frac{1}{p^4}\,,
  \qquad p^2/\Lambda^2 \to \infty\,.
  \label{eq:UV-decay}
\end{equation}
This is a consequence of the smoothness of $f$ and is specific to the
spectral action (the one-loop effective action has logarithmic growth
instead).


%%======================================================================
\newpage
\section*{Summary of Gaps}
\addcontentsline{toc}{section}{Summary of Gaps}
%%======================================================================

The following items are \textbf{not available} in the extracted
literature and must be computed in the NT-1 derivation:

\begin{enumerate}
\item \textbf{Explicit Dirac trace reduction:}
  The spinor trace of the five form factor terms in
  \eqref{eq:heat-trace-dirac-R2}, converting from the
  $\{R_{\mu\nu},R,\hat{P},\Omega_{\mu\nu}\}$ basis to the
  Weyl basis $\{C_{\mu\nu\rho\sigma},R\}$ with all Dirac-specific
  coefficients computed.

\item \textbf{The $s$-integration against the spectral weight:}
  The explicit integral \eqref{eq:F1F2-integral} converting
  heat kernel form factors $h_C$, $h_R$ into spectral action
  form factors $F_1$, $F_2$ for a general spectral function $f$.
  \emph{This is the core new computation.}

\item \textbf{BV parametric weights $\Phi_1$, $\Phi_2$, $\Phi_3$:}
  The exact polynomial weight functions in the parametric
  representation \eqref{eq:BV-parametric} for $f_1$, $f_2$, $f_3$.
  These are available in BV~(1987) Nucl.\ Phys.\ B \textbf{282}
  but were not extracted here.

\item \textbf{Closed-form result for specific $f$:}
  For common choices of spectral function (e.g., $f(u)=e^{-u}$,
  or $f(u)=\chi_{[0,1]}(u)$), the integrals in
  \eqref{eq:F1F2-integral} can be evaluated in closed form.
  These are not in the literature.
\end{enumerate}


%%======================================================================
\section*{References}
\addcontentsline{toc}{section}{References}
%%======================================================================

\begin{thebibliography}{99}

\bibitem{BV1985}
  A.~O.~Barvinsky and G.~A.~Vilkovisky,
  \emph{The generalized Schwinger--DeWitt technique in gauge theories
  and quantum gravity},
  Phys.\ Rept.\ \textbf{119} (1985) 1--74.

\bibitem{BV1987}
  A.~O.~Barvinsky and G.~A.~Vilkovisky,
  \emph{Beyond the Schwinger--DeWitt technique: converting loops into
  trees and in-in currents},
  Nucl.\ Phys.\ B \textbf{282} (1987) 163--188.

\bibitem{BV1990a}
  A.~O.~Barvinsky and G.~A.~Vilkovisky,
  \emph{Covariant perturbation theory (II). Second order in the
  curvature. General algorithms},
  Nucl.\ Phys.\ B \textbf{333} (1990) 471--511.

\bibitem{BV1990b}
  A.~O.~Barvinsky and G.~A.~Vilkovisky,
  \emph{Covariant perturbation theory (III). Spectral representations
  of the third-order form factors},
  Nucl.\ Phys.\ B \textbf{333} (1990) 512--524.

\bibitem{BGVZ2009}
  A.~O.~Barvinsky, Yu.~V.~Gusev, V.~V.~Zhytnikov, and
  G.~A.~Vilkovisky,
  \emph{Covariant perturbation theory (IV). Third order in the
  curvature},
  arXiv:0911.1168 [hep-th].

\bibitem{BGVZ1994}
  A.~O.~Barvinsky, Yu.~V.~Gusev, G.~A.~Vilkovisky, and
  V.~V.~Zhytnikov,
  \emph{Asymptotic behaviours of the heat kernel in covariant
  perturbation theory},
  gr-qc/9404063.

\bibitem{Barvinsky2014}
  A.~O.~Barvinsky,
  \emph{Aspects of nonlocality in quantum field theory, quantum gravity
  and cosmology},
  arXiv:1408.6112 [hep-th].

\bibitem{BW2023}
  A.~O.~Barvinsky and W.~Wachowski,
  \emph{Notes on conformal anomaly, nonlocal effective action and the
  metamorphosis of the running scale},
  arXiv:2306.03780 [hep-th].

\bibitem{Avramidi1990a}
  I.~G.~Avramidi,
  \emph{The nonlocal structure of the one-loop effective action via
  partial summation of the asymptotic expansion},
  Phys.\ Lett.\ B \textbf{236} (1990) 443--449.

\bibitem{Avramidi1990b}
  I.~G.~Avramidi,
  \emph{The covariant technique for calculation of the heat kernel
  asymptotic expansion},
  Phys.\ Lett.\ B \textbf{238} (1990) 92--97.

\bibitem{Avramidi1997}
  I.~G.~Avramidi,
  \emph{Covariant techniques for computation of the heat kernel},
  arXiv:hep-th/9704166.

\bibitem{Avramidi2000}
  I.~G.~Avramidi,
  \emph{Heat Kernel and Quantum Gravity},
  Lect.\ Notes Phys.\ Monogr.\ \textbf{64}, Springer (2000).

\bibitem{Vassilevich2003}
  D.~V.~Vassilevich,
  \emph{Heat kernel expansion: user's manual},
  Phys.\ Rept.\ \textbf{388} (2003) 279--360,
  arXiv:hep-th/0306138.

\bibitem{CZ2012}
  A.~Codello and O.~Zanusso,
  \emph{On the non-local heat kernel expansion},
  J.\ Math.\ Phys.\ \textbf{54} (2013) 013513,
  arXiv:1203.2034 [math-ph].

\bibitem{CPR2009}
  A.~Codello, R.~Percacci, and C.~Rahmede,
  \emph{Investigating the ultraviolet properties of gravity with a
  Wilsonian renormalization group equation},
  Ann.\ Phys.\ \textbf{324} (2009) 414--469,
  arXiv:0805.2909 [hep-th].

\bibitem{ILV2011}
  B.~Iochum, C.~Levy, and D.~Vassilevich,
  \emph{Spectral action beyond the weak-field approximation},
  Commun.\ Math.\ Phys.\ \textbf{316} (2012) 595--613,
  arXiv:1108.3749 [hep-th].

\bibitem{ILV2012}
  B.~Iochum, C.~Levy, and D.~Vassilevich,
  \emph{Global and local aspects of spectral actions},
  J.\ Phys.\ A \textbf{45} (2012) 374020,
  arXiv:1201.6637 [math-ph].

\bibitem{BMZ1995}
  A.~O.~Barvinsky, A.~G.~Mirzabekian, and V.~V.~Zhytnikov,
  \emph{Conformal decomposition of the effective action and covariant
  curvature expansion},
  gr-qc/9510037.

\bibitem{BG1995}
  A.~O.~Barvinsky and Yu.~V.~Gusev,
  \emph{Heat kernel and loop currents by the generating function method},
  gr-qc/9507026.

\bibitem{BGVZ1994b}
  A.~O.~Barvinsky, Yu.~V.~Gusev, G.~A.~Vilkovisky, and
  V.~V.~Zhytnikov,
  \emph{The one-loop effective action and trace anomaly in four
  dimensions},
  arXiv:hep-th/9404187.

\bibitem{CC1997}
  A.~H.~Chamseddine and A.~Connes,
  \emph{The spectral action principle},
  Commun.\ Math.\ Phys.\ \textbf{186} (1997) 731--750,
  arXiv:hep-th/9606001.

\end{thebibliography}

\end{document}
